\documentclass[11pt]{article}
\usepackage[a4paper,margin=25mm]{geometry}
\usepackage{amsmath,amssymb,amsthm}
\newtheorem{lemma}{Lemma}
\newcommand{\N}{\operatorname{N}}
\title{Incidence weights from a list with repeated representatives}
\author{}
\date{}
\begin{document}\maketitle

The descent and transport results in the accompanying manuscript assume a
complete list of distinct local conjugacy classes. A software interface may
return a complete list with repetitions. The following elementary observation
allows us to use that weaker input. It changes neither the descent theorem
nor the classification being computed. We claim no new group-theoretic result.

\begin{lemma}
Let $H$ be a finite group. Suppose $M_1,\ldots,M_s$ are maximal proper
subgroups of $H$, and every maximal subgroup of $H$ is $H$-conjugate to at
least one $M_i$. Process the list in order, retaining $M_i$ precisely when it
is not $H$-conjugate to any previously retained subgroup. The retained list
contains exactly one representative of each $H$-conjugacy class of maximal
subgroups. Its completeness does not require the original list to be
irredundant or to carry correct multiplicities.
\end{lemma}
\begin{proof}
Two retained subgroups cannot be conjugate: the later one would have been
discarded. For completeness, let $M$ be any maximal subgroup and choose an
index $i$ for which $M_i$ is conjugate to $M$. If $M_i$ is retained, its
class is represented. Otherwise the test that discards it supplies a
previously retained representative of the same class. Thus every class is
represented and no class is represented twice. All retained subgroups are
maximal because they belong to the original list.
\end{proof}

If $H$ is replaced by an isomorphic permutation group $\widetilde H$, this
operation is performed using conjugacy in $\widetilde H$. The transport
lemma carries the resulting distinct classes back to $H$. Conjugacy in a
larger ambient group must not be used for this step: distinct local classes
may fuse in that group, and each then contributes to the incidence number.

\begin{lemma}
Let $H,K\leq G$ be finite groups. Among the retained representatives above,
let $M_{i_1},\ldots,M_{i_t}$ be precisely those that are $G$-conjugate to $K$.
Then the number of actual maximal subgroups of $H$ in the ambient class
$[K]_G$ is
\[
 a_{H,K}=\sum_{j=1}^t[H:\N_H(M_{i_j})].
\]
Writing $c_H=[G:\N_G(H)]$ and $c_K=[G:\N_G(K)]$, the reverse incidence is
$b_{K,H}=c_Ha_{H,K}/c_K$. Necessarily
\[
 0\leq a_{H,K}\leq c_K,\qquad 0\leq b_{K,H}\leq c_H.
\]
In particular, if $K$ is maximal in $G$, then $a_{G,K}=c_K$ and
$b_{K,G}=1$.
\end{lemma}
\begin{proof}
Conjugation by $H$ has stabilizer $\N_H(M_{i_j})$ at $M_{i_j}$, so its orbit
has the displayed index. The first lemma makes these orbits pairwise
disjoint, and its completeness shows that their union is precisely the
set being counted. Counting incident pairs of ambient conjugates first
by their larger and then by their smaller subgroup gives the formula for
$b_{K,H}$. Each fibre is a subset of the appropriate ambient conjugacy
class, which proves both bounds. When the larger subgroup is $G$, all
conjugates of a maximal $K$ are maximal in it, and there is just one
ambient subgroup conjugate to $G$.
\end{proof}

The implementation therefore discards repeated local conjugacy classes
before ambient fusion and obtains each weight from a normalizer index.
It does not use a multiplicity attached to the original list. The two bounds
provide useful necessary checks; satisfying them alone does not establish
correctness or completeness. Completeness of the original maximal-subgroup
list and exactness of the group operations remain explicit computational
premises.
\end{document}
